\documentclass[12pt]{article}

\usepackage{sbc-template}
\usepackage{graphicx,url}
\usepackage[utf8]{inputenc}
\usepackage[brazilian]{babel}
\usepackage{etoolbox}

\sloppy
\setlength{\parskip}{1pt}
\setlength{\abovecaptionskip}{3pt}
\setlength{\belowcaptionskip}{0pt}
\setlength{\intextsep}{4pt}
\setlength{\textfloatsep}{6pt}

% Trabalho de Segurança em Redes de Computadores, limitado a quatro páginas.
% Formato SBC.
\title{Reconhecimento de Entidades em Incidentes Cibernéticos na Saúde: Reconstrução e Extensão de um Estudo}

\author{Gabriel S. Rodrigues\inst{1}, Gabriel Machado\inst{1}, Angelo Rodrigues\inst{1}
}

\address{Universidade Federal do Sul e Sudeste do Pará (UNIFESSPA) \\
 Marabá, Pará, Brasil
 \email{\{gabriel.s\_r, gabrielcarlosnm, angelo.rodrigues\}@unifesspa.edu.br}
}

\pdfinfo{
  /Title (Reconhecimento de Entidades em Incidentes Ciberneticos na Saude: Reconstrucao e Extensao de um Estudo)
  /Author (Gabriel S. Rodrigues, Gabriel Machado, Angelo Rodrigues)
  /Subject (Reconstrucao metodologica e extensao de uma avaliacao de NER)
}

\begin{document}

\maketitle

\begin{resumo}
  \normalfont
  Reconstruímos a avaliação das variantes NER e APTNER do SecureBERT em notícias sobre incidentes cibernéticos na saúde, proposta por Paim e Ignaczak. Como o corpus e as anotações originais não estavam disponíveis, coletamos 299 notícias e utilizamos anotações geradas por um modelo de linguagem, posteriormente revisadas manualmente. Preservamos as categorias organização, ação de ataque, vulnerabilidade e localização. Na referência revisada, os modelos obtiveram F1 estrito de 0,0809 e 0,0507, respectivamente, com baixa precisão e revocação. Como extensão, avaliamos o ajuste fino completo e a influência de uma anotação mais abrangente. Os resultados mostram sensibilidade à política de anotação e à regra de correspondência entre entidades. O estudo constitui uma reconstrução metodológica com dados próprios, não uma reprodução numérica do original.
  \textbf{Palavras-chaves:} reconhecimento de entidades nomeadas, cibersegurança, notícias de saúde, replicação.
\end{resumo}

\begin{abstract}
  \normalfont
  We reconstruct the evaluation of the NER and APTNER variants of SecureBERT on healthcare cybersecurity news proposed by Paim and Ignaczak. Since the original corpus and annotations were unavailable, we collected 299 articles and used annotations generated by a language model and subsequently reviewed manually. We retained the organization, attack action, vulnerability, and location categories. Against the revised reference, the models achieved strict F1 scores of 0.0809 and 0.0507, respectively, with low precision and recall. As an extension, we evaluated full fine tuning and the effect of more comprehensive annotation. Results show sensitivity to the annotation policy and entity matching rule. This study is a methodological reconstruction using a new corpus, not a numerical reproduction of the original.
  \textbf{Keywords:} named entity recognition, cybersecurity, healthcare news, replication.
\end{abstract}

\section{Introdução}

Incidentes cibernéticos que afetam hospitais e operadores de saúde têm impacto direto na continuidade do atendimento e na confidencialidade de dados de pacientes. Notícias especializadas registram esses eventos em linguagem natural e constituem uma fonte candidata para sistemas de extração de informações. Paim e Ignaczak \cite{paim2025} investigaram o uso de modelos pré-treinados \cite{devlin2019bert} para identificar entidades de incidentes cibernéticos em notícias da área da saúde e avaliaram SecureBERT-NER e SecureBERT-APTNER em uma amostra anotada manualmente. O corpus e as anotações originais, porém, não foram disponibilizados, o que impede a reprodução direta dos resultados.

Este artigo descreve a construção de um conjunto de dados próprio que preserva a tarefa e as quatro categorias do protocolo original, mantendo a distinção entre pré-anotação assistida e camada adjudicada. As contribuições são: (i)~um corpus de 299 notícias completas, documentado quanto à coleta; (ii)~duas camadas independentes, com 2.103 e 5.111 entidades; (iii)~avaliação sistemática de modelos e receitas de ajuste fino, com experimentos cruzados; e (iv)~uma ferramenta local de auditoria que permite categorizar entidades e adicionar trechos omitidos pela pré-anotação.

A seção~\ref{sec:corpus} descreve a construção do corpus e o protocolo de anotação. A seção~\ref{sec:exp} apresenta os experimentos e os resultados. A seção~\ref{sec:conc} discute limitações, disponibilidade e trabalhos futuros.

\section{Reconstrução do experimento}
\label{sec:corpus}

\subsection{Construção do corpus}

As URLs vieram do conjunto público \texttt{vinitvek/cybersecurityattacks}\footnote{\url{https://huggingface.co/datasets/vinitvek/cybersecurityattacks}}, filtrado por setor de saúde e pelo intervalo de 1 de janeiro a 30 de outubro de 2023. Os links serviram como índice para coletar, com Firecrawl, os textos completos nas páginas de origem, e não apenas os resumos do conjunto público. Das 329 URLs únicas selecionadas, 314 tiveram extração concluída. Após excluir 15 páginas inutilizáveis, restaram 299 notícias.

O original avaliou 60 notícias anotadas manualmente, limitadas aos primeiros 400 \textit{tokens}. Aqui, avaliamos textos completos com anotações automáticas revisadas manualmente. Portanto, mantivemos a tarefa, mas modificamos a seleção, a extensão dos textos e a produção das referências.

\subsection{Referências de anotação}

A anotação partiu de sugestões geradas pelo Muse Spark 1.3, um modelo de linguagem desenvolvido pela Meta \cite{meta2026muse}. Os anotadores revisaram manualmente os trechos e suas categorias, corrigiram marcações e acrescentaram entidades ausentes. Cada anotação delimita um trecho do texto e atribui uma categoria. Foram usadas quatro categorias. \texttt{IDTY} identifica organizações ou instituições. \texttt{ACT} corresponde a ações ou tipos de ataque. \texttt{VUL} marca vulnerabilidades. \texttt{LOC} identifica localidades geográficas.

A referência revisada, denominada \textit{overlay} nos artefatos, reúne 5.111 entidades nos 299 documentos. São 2.627 IDTY, 1.747 ACT, 65 VUL e 672 LOC. A referência de consenso contém 9.318 entidades e inclui menções repetidas da mesma entidade ao longo do texto. Na consolidação, foram mantidos os trechos da revisão final. Quando os mesmos limites apareciam nas outras anotações, os rótulos foram conciliados por voto majoritário. Em empate, prevaleceu o rótulo da revisão final. Trechos removidos ou corrigidos nessa revisão não foram reintroduzidos.

A cobertura não foi igual entre os anotadores. Dos trechos da referência de consenso, 4.640 têm três registros de anotação, 10 têm dois e 4.668 pertencem apenas à revisão final. Assim, as medidas de concordância entre os anotadores se referem aos trechos em comum, não à totalidade das 9.318 entidades.

Uma ferramenta local apoiou a conferência e a inclusão de entidades ausentes. Validadores automáticos verificaram rótulos e posições dos trechos, mas não garantem correção semântica. Nenhuma dessas referências equivale à anotação manual original.

\section{Experimentos e resultados}
\label{sec:exp}

\subsection{Protocolo de avaliação}

Foram avaliados os \textit{checkpoints} públicos SecureBERT-NER e SecureBERT-APTNER, com revisões fixadas e registradas no artefato de resultados, em inferência sobre os 299 documentos. Os artigos completos foram processados com limite de 512 \textit{tokens}, janelas alinhadas a palavras e \textit{stride} de 128. Para medir a influência da camada de anotação, a avaliação foi repetida sobre a camada Muse e sobre o \textit{overlay}. Três definições de correspondência foram empregadas: \textit{strict}, que exige limites e tipo exatos; \textit{partial}, que considera sobreposição de \textit{spans} sem exigir o tipo; e \textit{partial same-type}, que exige sobreposição e tipo correto. Essas métricas não são intercambiáveis.

Para o ajuste fino do SecureBERT-NER, adotou-se divisão documental fixa de 179 documentos para treino, 60 para validação e 60 para teste (\textit{seed}~2024), cinco épocas, sequência máxima de 512 \textit{tokens} e \textit{stride}~128. A cabeça BIO possui nove rótulos (\texttt{O} e \texttt{B/I} para as quatro categorias). Foram comparadas quatro receitas com três \textit{seeds}: atualização completa (\textit{Full}), \textit{adapters} de baixo \textit{rank} (\textit{LoRA+}), entropia cruzada ponderada (\textit{CCE}) e \textit{Repeated OverSampling} (\textit{ROS}). Resultados são médias~$\pm$~desvio padrão entre \textit{seeds}.

\subsection{Inferência e comparação entre camadas}

A Tabela~\ref{tab:layers} apresenta a comparação entre camadas. Na avaliação contra o \textit{overlay}, o F1 \textit{strict} do SecureBERT-NER passou de 0,0450 para 0,0809 e o F1 \textit{partial}, de 0,1531 para 0,2617. No \textit{overlay}, o SecureBERT-APTNER obteve F1 \textit{partial} 0,2933 e F1 \textit{partial same-type} 0,2891, maiores que o NER, mas menor F1 \textit{strict}.

Na camada aumentada, o SecureBERT-NER atingiu 0,0852 de F1 \textit{strict} e 0,2795 de F1 \textit{partial}, e o SecureBERT-APTNER, 0,0532 e 0,3111.

\begin{table}[ht]
\centering
\caption{Inferência na referência revisada (299 notícias). P e R são precisão e revocação no critério estrito.}
\label{tab:inferencia}
\begin{tabular}{l|c|c|c|c}
\hline
\textbf{Modelo} & \textbf{P} & \textbf{R} & \textbf{F1 estrito} & \textbf{F1 parcial} \\
\hline
SecureBERT NER & 0,0896 & 0,0738 & 0,0809 & 0,2617 \\
SecureBERT APTNER & 0,0622 & 0,0428 & 0,0507 & 0,2933 \\
\hline
\end{tabular}
\end{table}


\subsection{Ajuste fino}

Na camada \textit{overlay}, \textit{Full}, \textit{Full+CCE} e \textit{Full+ROS} atingiram 0,4336~$\pm$~0,0056, 0,4421~$\pm$~0,0049 e 0,4410~$\pm$~0,0073 de F1 \textit{strict}, contra 0,1061~$\pm$~0,0152 do \textit{LoRA+}. Em F1 \textit{partial same-type}, o \textit{Full+ROS} obteve 0,5602~$\pm$~0,0038. Experimentos cruzados com o \textit{Full} reforçam a sensibilidade à camada: treino Muse/teste \textit{overlay} atingiu 0,2403 de F1 \textit{strict} e 0,3699 de F1 \textit{partial same-type}; a configuração inversa, 0,2379 e 0,3033; a diagonal \textit{overlay}/\textit{overlay}, 0,4336 e 0,5528.

No conjunto de teste, uma \textit{seed} preliminar (42) na camada aumentada atingiu 0,4866/0,6344 de F1 \textit{strict}/\textit{partial same-type}, contra 0,3058/0,4558 na Muse na mesma \textit{seed}. O resultado é preliminar e não substitui a média de três \textit{seeds}.

\section{Limitações e conclusão}
\label{sec:conc}

O conjunto tem três limitações: o \textit{overlay} é uma camada adjudicada, não o \textit{gold} manual; os textos e anotações por documento são restritos; e o ajuste fino é específico da divisão e das receitas testadas.

Em síntese, o recurso reúne 299 notícias, duas camadas de anotação e 687 adições assistidas por IA e verificadas por humano; a camada e a métrica devem ser declaradas.


{\normalsize
\vspace{4pt}
\noindent\textbf{Uso de IA.} O Muse Spark foi utilizado para gerar as anotações iniciais. Ferramentas de IA também apoiaram a redação e o desenvolvimento do código. A revisão das anotações e as decisões finais foram realizadas pelos anotadores humanos. Os autores respondem pelo conteúdo e pelas decisões científicas.

\bibliographystyle{sbc}
\bibliography{references}
}

\end{document}
